\section{Experiments}
\label{sec:experiments}
\subsection{Datasets, metrics and comparison protocol}
We use one fixed in-house method recipe across 31 scene jobs. Four LERF scenes yield
two outputs, so the cohort comprises five evaluation tasks rather than five
independent datasets. All reported numbers are percentages. The existing
benchmark rooms have been used during development; they are not an untouched
final test. Scene-level intervals describe this cohort rather than certify
generalization after repeated research feedback.

\begin{table}[ht]
\centering\small
\begin{tabular}{lp{.25\linewidth}p{.43\linewidth}}
\toprule
Task & Cohort and primary metric & Input and evaluation contract \\
\midrule
LERF-2D & 4 scenes, 208 observations; mIoU and localization & Source-assisted text; semantic response rendered before thresholding. \\
LERF-3D selection & Same scenes; projection mIoU and Acc.@0.25 & Independent selected Gaussians rendered to 2D GT; not volumetric GT IoU. \\
ScanNet OVS & Paper8 Gaussian mIoU/mAcc; internal10 mesh mIoU & Fixed 19/15/10 menus; Gaussian and mesh domains reported separately. \\
NVOS & 8 tasks, 7 scene clusters; IoU & Retained positive-reference, cached-source variant; external prompt/RGB access declared. \\
UQIS & 9 rooms, 67 instances, 268 queries; set IoU and AP & Text, image, 2D point and 3D point; map-only query-time access. \\
\bottomrule
\end{tabular}
\caption{Evaluation units. LERF and UQIS average within scenes then across
scenes; NVOS averages eight tasks, with three seeds per task for LUDVIG.}
\label{tab:protocols}
\end{table}

\paragraph{Published results and verified reproductions.}
We reuse values reported in VALA Tables 1--2 and LUDVIG Tables 1/D, including
their cited baseline results~\cite{vala,ludvig}. No new external baseline training
is needed. Every local reproduction is linked to its frozen result hashes;
four-scene, eight-scene and eight-task aggregates are recomputed from the
original reports. Table origins are \textbf{P}: published, \textbf{R}: locally
reproduced, \textbf{O}: our implemented method (v10 or a named internal control), and \textbf{M}: our method with a named measured extension.
Missing extension results remain pending in the registry; they do not erase the evaluated v10 method results.
We do not bold a global winner across different input contracts.

\subsection{Implementation and supervision}
The fixed v10 recipe includes native teacher matching, source evidence, fixed
geometry and typed task readouts. SAM3~\cite{sam3} source masks and native
SigLIP2~\cite{siglip2} scores support LERF. ScanNet combines original SigLIP2
and Mask-Adapter~\cite{maskadapter} priors with
literal-category SAM3 evidence. NVOS uses positive-prompt SAM3 reference tracks
and a robust rendered response. All nine UQIS maps use SigLIP2 giant and
DINOv2 ViT-B/14~\cite{dinov2} for image matching;
complete membership inherits supervision on 22 labelled rooms and calibration
on three other rooms. These room identities are disjoint from benchmark rooms.
The original feature fields and teacher weights are not retrained in this study.

The optional L11 selector additionally uses independent text/image/category
utility supervision. Its 22 fit, three calibration and four development rooms
provide 629/207/175 text expressions, 237/28/25 image queries and 116 category
episodes. Only 88 fit-room category episodes enter training; calibration and
development use full maps. Query-image views and neighbors are excluded from
their candidate observations; full-scan geometry remains available. One margin
(0.15052) protects 15 calibration room/output cells. Five room folds fit a
64-tree structural prior plus a 64-tree query residual. All 19 ScanNet/UQIS jobs
are predicted before scoring. These are extension settings, not prerequisites
for the v10 main results. Development and calibration informed model design.

\subsection{Comparison with existing methods}
\paragraph{LERF semantic segmentation and projected selection.}
Tables~\ref{tab:lerf2d}--\ref{tab:lerf3d} reuse the published context and verified
Occam/VALA reproductions. The VALA reproduction follows its semantic and
evaluation path on compatible Occam RGB geometry, rather than retraining the
entire original geometry pipeline. The v10 values are our method results. The optional shared replacement has
not yet been measured on these two outputs. In particular, semantic-response improvements cannot be substituted
for independent selected-set improvements.
Our method includes SAM3 text-conditioned source evidence and native SigLIP2
matching; external VALA scores are not matched-teacher controls. An independent
CPU audit exactly reproduces all 208 retained selected-set mask IoUs. Semantic
IoU/localization records are reaggregated, without a fresh semantic render in
that audit.

\begin{table}[ht]
\centering\small\input{tables/comparison_lerf2d.tex}
\caption{LERF-2D. P values are reported in VALA Table 1~\cite{vala}.
Localization is peak-in-box accuracy, not ranking AP. O uses source-assisted
full teachers; these are v10 method results, not new replacement gains.}
\label{tab:lerf2d}
\end{table}
\begin{table}[ht]
\centering\small\input{tables/comparison_lerf3d.tex}
\caption{LERF selected-Gaussian projection. P values are reported in VALA
Table 1~\cite{vala}; Acc.@0.25 counts projected masks with IoU above 0.25.
VALA's published value is 58.02/82.85; 43.29/64.30 belongs to Dr. Splat.
The local VALA R row uses compatible RGB geometry.}
\label{tab:lerf3d}
\end{table}

\paragraph{ScanNet open-vocabulary segmentation.}
Table~\ref{tab:scannet} separates the paper-eight Gaussian metric from the
original ten-room mesh metric. For the metric bridge, we score unchanged v10
labels with the existing VALA evaluator: anisotropic Mahalanobis pseudo labels,
five-scale radius, top-1000 class-balanced votes and opacity--volume weights.
All eight scene predictions are hashed before opening their target labels.
This aligns the metric and scene list, but not the input carrier: v10 uses
provided unlabelled mesh geometry, while VALA uses optimized RGB Gaussians.
The bridge is a protocol diagnostic, not an additional learned gain. The
replacement policy retains every original menu on the internal ten-room track.
Independent scoring reproduces all ten mesh scenes and all eight Gaussian
scenes. All 1,134,207 pseudo labels agree, including 256 checks with the released
VALA code. On the same eight rooms, 19-class mesh mIoU is 61.0567 versus
66.9480 for Gaussian scoring: this 5.8913-point change is metric-dependent.

\begin{table}[ht]
\centering\small
\resizebox{\linewidth}{!}{\input{tables/comparison_scannet.tex}}
\caption{ScanNet OVS. P rows from VALA Table 2~\cite{vala}, R and the metric
bridge use paper8 Gaussian scoring. Rows marked 10-mesh use the original
ten-room mesh evaluation and cannot be ranked against the paper8 rows.
The metric bridge still differs in geometry and teacher access.}
\label{tab:scannet}
\end{table}

\paragraph{NVOS.}
Table~\ref{tab:nvos} includes the complete three-seed LUDVIG-SAM reproduction.
Its positive scribbles, all-view RGB geometry and online target-RGB access
follow the released implementation. Our original v10 method uses a
positive-derived reference SAM mask and cached source tracks, without online
target RGB. These are contextual method comparisons with declared access
differences. Later image-conditioned extensions explicitly permit observed
target RGB; they retain the original geometry and are reported separately.
The audit also finds a GT boundary convention difference: our retained scorer
uses $g>127$, while LUDVIG uses integer division by the mask maximum. Applying
only its GT convention changes our score from 86.5134 to 86.6835; historical
values remain unchanged. Switching bilinear implementations changes less than
0.00002 points. Borrowing its min--max and fixed 0.25 response threshold yields
86.6085 on our GT convention, and is not a full LUDVIG reproduction. Flower,
trex and orchids account for 78.04\% of the gap to the local 91.2577 reproduction.
Flower/orchids primarily miss foreground; trex also activates background.
N04 pairs reconstructed and observed RGB at the same camera, with identical
response-derived prompts and frozen image decoder. Observed RGB increases
IoU from 87.5575 to 89.0636 for direct image responses, and from 86.9242 to
88.3201 after actual 3D writeback. This isolates an image-access intervention
within this decoder; it does not explain the entire difference to LUDVIG.
The 0.7435-point loss between direct image and 3D outputs leaves transport,
observation pooling, geometry and the fitting objective to be separated.
The 3D candidate raises IoU by 1.8067 points over v10 but lowers AP, so it is
not a uniform improvement across metrics. Source-teacher residuals also rise;
this is a view-conditioned update, not a demonstrated query-global field.
N05 finds that discarding negative dense hints loses 2.6331 IoU points against
the matched signed-hint control. Sparse hints recover flower but admit another
object in horns-center. Complete prompt controls and both GT conventions are
reported in the diagnostic supplement; no per-task winners are combined.
N07's reference-memory transfer does not exceed the static image-conditioned
mean. N08 fixes all three image-evidence arrays and fits responses on the native
1008-by-756 ray grid instead of 512-by-384. The predeclared sequence-memory
primary arm reaches 88.1124 IoU. The static secondary arm reaches 88.8424,
up 0.5223 points over its coarse control with all eight tasks improved, leaving
a 0.2212-point gap to its unchanged direct image evidence. An accompanying
CUDA arithmetic correction changes matched coarse-grid mean IoU by less than
0.000006 points. These results support resolution-consistent transport, not
a benefit from memory itself. N09 replays this static control with identical
per-task binary counts in both GT tracks and measures 96.5817 AP, below v10's
98.0222. The static image arrays match exactly; the native continuous arrays
are not bitwise identical to N08. Neither run certifies SOTA.

N09 adds frozen DINOv2 descriptors and the official negative reference scribbles
to test complete query transfer beyond the initial mask. Signed-scribble authority
and reference-mask/complement authority use matched inputs and downstream
operations. Their native IoUs are 81.5803 and 86.0812. Signed evidence recovers
missing regions in flower and orchids, but confuses other instances in
horns-center and leaves. Its source leave-one-out accuracy of 99.80--100\%
therefore does not establish target-instance correctness. N10 writes the exact
affinity probabilities directly to 3D, yielding 77.7307 and 84.2098 IoU;
omitting SAM does not recover the static mean. All 72 N09 and 48 N10 predictions
are sealed before their respective scoring. Detailed paired controls, AP and
full-cohort error maps appear in the diagnostic supplement. The added teacher
and negative-prompt access are disclosed; no per-task branch selection is used.

N11 fixes those image proposals and compares unrestricted additive completion,
v10-response anchoring and signed-scribble geometry anchoring. Their actual 3D
IoUs are 83.5805, 89.9673 and 89.9734. The predeclared geometry primary gains
1.1310 points over static native writeback (six tasks improve, two regress),
but loses 0.1540 AP points. Flower improves from 85.9782 to 94.9302, while
instance constraints prevent the unrestricted branch's severe horns-center
and leaves regressions. Orchids remains 73.5362: the rule rejects most of the
proposal's missing-flower recovery. The 0.0061-point difference between the
two anchor policies does not establish superiority of geometric witnesses.
All 72 predictions precede scoring. N11 is a development improvement, not a
SOTA certification or a common improvement across all five benchmarks.
N12 moves the intervention to the reference image. Original, semantic-completed
and signed-constrained reference masks enter identical two-frame memory and
native fitting, giving 88.2002, 88.5226 and 91.4774 IoU. Constrained completion
raises orchids to 94.5224 but harms horns-left by modifying a reference mask
that already agrees with all original scribbles. N13 instead transfers the
paired memory difference around the static response: probability and log-odds
variants give 90.7469 and 89.4416. The probability residual harms all eight tasks
versus its matched complete replacement, so it is not retained as a required
method stage.

N14 fixes two source-only admission rules after this N12 development diagnosis.
Feasibility admission (primary) reaches 92.9424 IoU / 98.8599 AP; source-quality
admission (secondary) reaches 93.1730 / 98.9451. Their LUDVIG-GT IoUs are 92.9852
and 93.2104. Quality exceeds static native writeback by 4.3305 points and N11
by 3.1996, improving five and harming three tasks in either comparison. It
improves only trex against feasibility, while harming fern and horns-center;
five tasks tie. Thus the aggregate does not establish a universally better
quality criterion. Flower, orchids and horns-left reach 96.7139, 95.6011 and
94.8079. The quality image output is 93.4971, with 0.3241 mean loss after actual
3D fitting; internal background leakage and fine structures remain limiting.

All 72/56/56 N12/N13/N14 outputs respectively precede their GT scoring. N14
freshly replays selected teacher responses and verifies exact parent hashes,
then fits and re-scores actual 3D outputs; it does not substitute cached
per-task maxima. Original negative scribbles, DINOv2 and target RGB are shared
with N11 but extend the original v10 input contract. The other 23 jobs remain
unchanged references. This is adaptive development evidence, without a matched
external SOTA certification or all-five default promotion.

N15 freezes the N14 reference-quality decision and transports both original and
completed references to all 318 authorized source views. Matched source-only
fusion reaches 88.0061 and 92.7461 IoU, respectively: completion gains 4.7400
points with no target-image evidence in either field. Equal-group fusion with
the N14 target response reaches 93.0302 IoU / 98.9797 AP, versus the replayed
N14 control's 93.1730 / 98.9451. It reduces false positives but also true
foreground on every task; precision increases while recall falls. Both GT
conventions retain the same mean-IoU ordering. The predeclared joint candidate
therefore fails the higher-IoU, non-decreasing-AP criterion, and this fusion
branch is stopped. All 48 outputs precede scoring, and the three N14/v10
controls reproduce the original binary counts. The source-only result supports
transfer of reference-completion capability; it does not establish that more
source constraints improve the strongest target-conditioned output.

N16 re-decodes the exact complete N14 target response once, using either the
full target image or a response-derived crop with 25 percent context per side.
The same crop transforms RGB and prior; outside probabilities are retained
before native fitting. Full-frame decoding gives 93.0998 IoU / 98.2626 AP.
The local primary reaches 93.2805 / 98.4663, versus the replayed N14 control's
93.1730 / 98.9450. Thus the highest observed native IoU increases by 0.1075
points, while AP loses 0.4787. Four tasks improve and four regress. Local
decoding exceeds its matched full-frame control by 0.1806 IoU and 0.2037 AP
points, but does not retain the N14 ranking quality. Fern and horns-center
receive full-frame windows under the fixed rule and have identical predictions
in the two controls. Both GT conventions show the tradeoff; omitting the whole
horns scene removes the mean-IoU gain over N14. All 56 predictions precede
scoring. The branch fails its predeclared joint-metric criterion and is stopped;
N16 remains a measured IoU/AP tradeoff, while N14 remains the retained balanced
recipe. Neither result certifies an external same-protocol SOTA ranking.
An exact eight-task score audit (A24, diagnostics) preserves all hard masks and
changes mean image AP by less than 0.002 points under higher-precision probability
interpolation; it does not overturn the N16 stopping decision or alter 3D scores.

\begin{table}[ht]
\centering\small\resizebox{\linewidth}{!}{\input{tables/comparison_nvos.tex}}
\caption{NVOS IoU. Our rows use all eight tasks. P denotes published values from
LUDVIG Tables 1/D~\cite{ludvig}, VCAR Table 1~\cite{vcar}, LEGO Table 1~\cite{lego}, WildSeg3D Table 1~\cite{wildseg3d}, and Trace3D Table 4~\cite{trace3d}.
VCAR uses seven targets; LEGO reports seven scenes, without a verified matching
target roster. These recent rows are contextual, not same-protocol reproductions.
R averages all three seeds in each task. O uses our declared reference/source
contract. New NVOS evidence variants are reported as controlled extensions.}
\label{tab:nvos}
\end{table}

\paragraph{Multimodal UQIS.}
UQIS is the local ScanNet-UQIS9 v0.3.1 protocol; no external published row in
the cited papers evaluates these same 268 queries. Table~\ref{tab:uqis} therefore
reports all modalities and a retained regional-readout baseline. The difference
between regional readout and v10 measures our complete-decoder capability;
the optional selector has its own separate comparison. The matched v10-to-replacement comparison isolates the
response decision. The full
shared policy improves image IoU by 1.088 points and AP by 1.268 points. One
query improves and none are harmed. Text and both actual-point outputs remain
unchanged; the point rows verify preservation rather than learned improvements.
All reported outputs use the same policy, without choosing a different winning
variant for text, images or categories.
\begin{table}[ht]
\centering\small\resizebox{\linewidth}{!}{\input{tables/comparison_uqis.tex}}
\caption{UQIS matched comparison. T/I/P2/P3 denote text, image, 2D point and 3D
point. Actual selected-set IoU and continuous-response AP are distinct metrics.}
\label{tab:uqis}
\end{table}

\subsection{Historical component controls within our method}
Tables~\ref{tab:core-components}--\ref{tab:extent-components} connect v10's
capability to actual interventions. They are historical experiments on the
stated cohorts, not newly invented leave-one-out scores of the final system.

\begin{table}[ht]
\centering\small
\resizebox{\linewidth}{!}{\input{tables/mechanism_evidence.tex}}
\caption{Our matched identity, response-use and image-extent controls.
First two rows use LERF selected projection and semantic IoU respectively;
last two use actual UQIS image instance IoU. Different controls are non-additive.}
\label{tab:core-components}
\end{table}
Query admission changes candidate rejection with geometry, features and the set
solver fixed. Response-use control holds evidence, map and independent set
fixed while changing how admitted observations enter the semantic output; its
19.79-point gain is not attributed to a newly invented averaging equation.
Image local correspondence changes matching with the same teachers and extent
head. Complete-carrier reuse changes complete membership from the same identity
condition. These controls support different links of the same response process.

\begin{table}[ht]
\centering\small
\resizebox{\linewidth}{!}{\input{tables/provider_evidence.tex}}
\caption{Additional in-house component controls. Dual-native fusion uses three
held-out ScanNet rooms; point rows use UQIS actual-set IoU. Point supervision
changes observation distribution and sample counts, not only a formula.}
\label{tab:extent-components}
\end{table}
The 2.35-point category gain includes a second native teacher's capability.
Real-observation membership controls change training-observation distribution
and counts under the same learner. They must not be attributed solely to an
evidence-weighting rule or reported as zero-supervision results.

Four later full-cohort controls preserve these capabilities while testing
partial/complement permissions, available-domain normalization, scope agreement
and response solvers. No candidate produces common gains across all scored
outputs. Squared reconstruction reduces source loss in all 22 source-assisted
jobs, but downstream quality is mixed. These negative controls support keeping
effective typed operators; they are detailed in the diagnostic supplement.

\IfFileExists{tables/nvos_coverage.tex}{\input{sections/nvos_coverage.tex}}{}
\IfFileExists{tables/nvos_source_readout.tex}{\input{sections/nvos_source_readout.tex}}{}
\IfFileExists{tables/nvos_local_credit.tex}{\input{sections/nvos_local_credit_summary.tex}}{}

\subsection{Optional complete-response replacement controls}

All rows in Table~\ref{tab:ablations} retain the same teachers, complete
candidates, supervision rooms and output definitions. The query-blind arm returns only the fitted structural prior (64 trees,
versus 128 in the full model); it therefore also changes capacity. Permutation
instead uses the fitted full model and
shuffles all teacher scores jointly within each view. It preserves their
multisets and correlations. Removing the null constraint retains the same
64+64-tree capacity and total fit weight, replacing permuted training examples
with duplicates of aligned examples. Each trained control uses the same
calibration procedure. The zero-margin control uses the full fitted model and
changes only the acceptance margin. These are response-level ablations.

\begin{table}[ht]
\centering\small\resizebox{\linewidth}{!}{\input{tables/mechanism_ablations.tex}}
\caption{Matched controls on all ten ScanNet rooms and 268 UQIS queries.
S19/15/10 are mesh mIoU. Repair/harm counts concern actual UQIS set IoU;
unchanged actual points are included in the denominator. This is one coupled
policy comparison, not a collection of per-benchmark best runs.}
\label{tab:ablations}
\end{table}

The image gain over query-blind and permuted controls is respectively 0.725 and
1.088 points; descriptive nine-room bootstrap intervals are [0,2.174] and
[0,3.264]. They do not establish a robust query-specific advantage. The full
and no-null arms make identical choices in all 144 text/image/menu decisions,
so the null constraint has no demonstrated accuracy contribution. Removing the
positive margin improves text and image IoU, but reduces ScanNet mIoU, including
a 0.869-point loss on the ten-class menu. Retention prevents some damage while
also rejecting useful instance replacements. Neither the full policy nor the
zero-margin variant is promoted to replace the reference system.

\subsection{Candidate capacity and failure attribution}
For a fixed candidate pool, let $G_0$ be gain under zero-margin selection and
$G_\tau$ gain under calibrated retention. The exact diagnostic is
\[
G_{\rm pool}=R_{\rm rank}+(G_0-G_\tau)+G_\tau.
\]
The middle term is signed: it is negative when retention prevents damage.
This identity prevents attributing all unused candidate capacity to the margin.
For text, 15.395 points of available IoU gain decompose into 9.615 ranking regret,
5.780 retention cost and zero actual gain. For images, 9.492 available points
decompose into 8.224 ranking regret, 0.180 retention cost and 1.088 actual gain.
Thus image ranking is the dominant measured bottleneck; lowering the margin
cannot recover most of the unused capacity. Category mean-menu capacity is
only 0.913 points, and its negative retention cost indicates prevented damage.
These label-oracle diagnostics are computed after prediction and are never
automatic method results.

\paragraph{Five-task coverage and remaining evidence.}
The implemented v10 method covers every task. Positive component controls cover
LERF identity/response use, ScanNet native priors and UQIS identity/extent;
NVOS requires evidence and geometry attribution beyond scorer equivalence.
The optional selector has separate missing LERF/NVOS transfer cells. These two
kinds of completion must not be conflated. Independent teacher configurations
and unseen scenes remain needed for broader claims; no feature-dimension or
compression sweep is required for this response-focused paper.
